\specialchapter{引言}

在本书第 III 章至第 V 章中，我们详细考察实或复拓扑向量空间的自同态代数。在有限维空间的情形，这一代数的元素逐个来看的主要性质都可由若尔当分解（A, VII, p. 42, § 5, 第 9 目）导出。我们在此研究它们向无穷维空间的推广，然而这推广要更为精细。

在第 III 章中，我们在非常一般的框架下考察拓扑向量空间之间的紧线性映射。它们的谱性质与有限维情形最为接近。例如，在巴拿赫空间的特殊情形，紧自同态的谱中的非零元素都是具有有限谱重数的广义特征值。

在希尔伯特空间的框架下，会出现一种易于处理的分类和结构。这正是第 IV 章的对象。该章从希尔伯特空间中紧算子的研究开始：它们借助奇异值的描述，以及正规自同态情形下的对角化，都完全类似于有限维空间上已知的性质。更一般地，谱定理给出了希尔伯特空间正规自同态的简单模型。

某些重要应用，例如涉及微分算子的应用，但也涉及量子力学的应用，要求一种更精微的理论：即定义在希尔伯特空间的稠密子空间上的线性映射---未必连续---的理论，我们把这类映射称为部分算子。在这一框架下，自伴算子或正规算子的定义本身就必须极其精确，否则便会产生重大的谬误。一旦掌握了这些定义，就会出现一种与连续自同态情形相类似的分类；而对有界或无界的普遍可测函数的函数演算，则在连续正规自同态与正规部分算子之间建立起紧密的联系。

顺带地，我们还会谈到带迹的正规自同态，特别是希尔伯特空间 \(L^{2}(X,\mu)\) 中带积分核的积分自同态。人们熟知，在自守形式与表示的理论中，断言这类自同态（当其存在时）的迹等于核沿对角线的积分的那条公式具有至关重要的意义。我们还用一节的篇幅叙述 \(R^{n}\) 中广义函数与缓增广义函数的基本定义和性质。这些概念不仅给出了欧几里得傅里叶理论诸基本原理（参见 II, § 1）的最方便的表述，而且还使人们能够确定许多重要微分算子的定义域。

最后，在第 V 章中，我们把希尔伯特空间自同态的研究与群的酉表示联系起来。在若干初等的一般性考虑之后---其中舒尔引理是最重要的命题---我们研究局部紧群的情形。然后我们考察正定函数，并给出它们的最初应用。这些应用之一是盖尔范德--奈马克--塞加尔定理，它断言每个星代数都同构于某个希尔伯特空间的自同态代数的一个闭子代数。另一些应用则涉及局部紧群的不可约酉表示所成的空间的基本性质，特别是盖尔范德--赖科夫定理，它确立了这些表示分离点。在本章最后一节中，我们证明关于紧群表示的基本结果，它们曾在《群与李代数》第 IX 章中为我们所用。

A. BOREL, Representations of locally compact groups, Lecture Notes in Mathematics 276, Springer, 1972.

M. REED 与 B. SIMON, Methods of Mathematical Physics, Academic Press, Vol. 1, 1980 与 Vol. 2, 1975.

K. SCHMÜDGEN, Unbounded Self-adjoint Operators on Hilbert Space, Graduate Texts in Mathematics 265, Springer, 2012.

在前面的各章与各书中，对“TS, to appear”的下列引用应作如下修改：

EVT, V, p. 35, 第 1 行。应将“TS, IV, § 6”改为“(TS, IV, p. 160, remark)”。

INT, IX, p. 94, § 6, 第 12 目。应将“One can prove a converse \ldots{} bounded”改为“The converse of prop. 11 is a particular case of TS, V, p. 455, th. 5.”

LIE, IX, p. 61, § 6, 第 4 目, 第 6 行。应将“TS, to appear”改为“TS, III, p. 36, prop. 9”。

LIE, IX, p. 63, § 6, 第 4 目, 第 6 行。应将“TS, to appear”改为“TS, III, p. 36, rem.”。

LIE, IX, p. 66, § 7, 注 1。应将“to a chapter to appear”改为“to paragraph 4 of Chapter V”。

LIE, IX, p. 69, § 7, 第 2 目。应将“then according to TS”改为“then (TS, V, p. 464, prop. 6)”。

LIE, IX, p. 71, § 7, 第 3 目。应将“which, according to TS, is an isomorphism”改为“which, according to TS, V, p. 469, cor. 2, is an isomorphism.”

LIE, IX, p. 72, § 7, 第 4 目。应将“(TS)”改为“(TS, V, p. 459, cor. 3)”。

LIE, IX, p. 74, § 7, 第 4 目。在“of the space of central representative functions on G”之后加上“(TS, V, p. 469, cor. 2, a)”；在“of the space ZL²(G) of central square-integrable functions on G”之后加上“(TS, V, p. 468, prop. 9).”

LIE, IX, p. 78, § 8, 第 1 目。应将“In this number, one recalls definitions and results from TS”改为“One recalls definitions and results from TS, V, p. 386, no. 8; p. 392, no. 10; p. 470, no. 7.”

LIE, IX, p. 84, § 8, 第 3 目。应将“Recall (TS) the formulas”改为“One has (TS, V, p. 467, formulas (1) and (2)) the formulas”。

LIE, IX, p. 89, § 9, 第 1 目。应将“TS, III, §2, no. 7, prop. 16”改为“TS, III, p. 67, prop. 10”。

LIE, IX, p. 90, § 9, 第 1 目。应将“TS, III, §2, no. 6”改为“TS, III, p. 61, no. 4”。

LIE, IX, p. 91, § 9, 第 2 目。应将“TS, to appear”改为“TS, III, p. 36, prop. 9 and TS, V, p. 464, prop. 7”。

LIE, IX, p. 92, 注 1。应将“(TS, to appear)”改为“(TS, V, p. 379, def. 3)”。

LIE, IX, p. 99, 附录 1, 第 1 目。应将“(TS, to appear)”改为“(TS, V, p. 464, remark 4)”。

LIE, IX, p. 128, § 9 习题 3 (b)。应将“(TS, to appear)”改为“(Use the Peter--Weyl theorem, TS, V, p. 462, th. 1)”。

TS, I, p. 103, 例 6。应将“V, to appear”改为“V, p. 442, th. 2.”

TS, I, p. 173, 习题 5, d)。应将“III, to appear”改为“III, p. 82, § 6.”

TS, II, p. 237, 第 9 目。应将“IV, to appear”改为“IV, p. 196, § 3”。

TS, II, p. 251, 第 1 目。应将“V, to appear”改为“V, p. 425, prop. 10”。

